This chapter presents \textbf{McpRisk}, our defense-oriented risk-scoring
framework. We first make the threat model precise
(\Cref{sec:framework:threat_model}), then describe the end-to-end system
architecture (\Cref{sec:framework:architecture}), the factor-based scoring
model and its formulas (\Cref{sec:framework:model}), the atomic-operation
severity taxonomy (\Cref{sec:framework:atomic}), the parameter layer that turns
the static score into a request-time score (\Cref{sec:framework:params}), and
finally the implementation and its reproducibility guarantees
(\Cref{sec:framework:implementation}).

\section{Threat Model}
\label{sec:framework:threat_model}
\addcontentsline{tocheb}{section}{\protect\numberline{\secnumforhebrewtoc}{מודל האיום}}

McpRisk adopts a single, deliberately asymmetric threat model:

\begin{quote}
\textbf{The MCP server is always the protected asset; the agent is always the
threat source.}
\end{quote}

Every interaction we consider flows \emph{client} $\to$ \emph{server}: an
autonomous or semi-autonomous LLM agent (or the user behind it) issues tool
calls through the protocol in order to read, modify, destroy, or exfiltrate
resources that the server fronts. The agent may be malicious by intent, or
benign but compromised --- for instance, an agent that has ingested an indirect
prompt injection~\cite{zhan2024injecagent} and now acts against the server's
interest. From the server's vantage point these cases are indistinguishable,
so McpRisk does not attempt to infer intent; it scores the \emph{consequence}
of the requested operation.

The inverse scenario --- a malicious server attacking its client, the classical
``tool poisoning'' direction~\cite{shi2025toolhijacker} --- is explicitly out
of scope. Tool poisoning enters our model only as a \emph{cause}: a poisoned
agent subsequently becomes a threat \emph{to the server}, and it is that
downstream request we score. Two further design commitments follow from this
framing. First, McpRisk is \textbf{not anomaly detection}: it does not flag
deviations from a learned baseline of ``normal'' traffic, because a
first-of-its-kind malicious call has no baseline. Instead it defines, from
concrete factors, what \emph{makes} an interaction risky. Second, the score is
\textbf{graduated, not binary}: rather than an allow/deny verdict, it produces a
four-level band so that a server can gate the top band, throttle or escalate the
middle bands, and log the rest.

\section{System Architecture}
\label{sec:framework:architecture}
\addcontentsline{tocheb}{section}{\protect\numberline{\secnumforhebrewtoc}{ארכיטקטורת המערכת}}

McpRisk is organized as a pipeline of loosely coupled stages, summarized in
\Cref{fig:pipeline}. It operates in two phases that mirror the two scoring
modes of the introduction.

\textbf{Design time (static).} A \emph{scanner} runs before the server goes
live. It never connects to a running server and reads no pre-existing risk
table; instead it assembles a registry from two static sources --- the tools
the server advertises in its own \texttt{tools/list} response, and the assets
of the store it fronts (a filesystem tree, a SQLite schema, or a set of
declared Slack channels, calendars, or repositories). A six-stage
\emph{static-scoring} pipeline then assigns each (tool, asset) pair a risk band,
producing a scan artifact.

\textbf{Request time (dynamic).} A \emph{parameter-scoring} module derives, per
tool, a rubric for the magnitude-bearing arguments of a call, and a
\emph{call-scoring} module resolves each captured call to a scanned (tool,
asset) cell and escalates its band according to the actual argument magnitudes.
Calls that cannot be resolved are handled honestly: an unknown tool is marked
\emph{invalid} (a misconfiguration) and a known tool aimed at an unscanned
target is marked \emph{unresolved} (the scorer abstains) rather than guessing.

A separate \emph{review} subsystem --- comprising deterministic verification, an
independent LLM judge, an advisor, and a panel of persona critics --- audits
the pipeline's outputs but is never imported by the scorer itself, keeping
evaluation independent of scoring.

\begin{figure}[H]
    \centering
    \begin{tikzpicture}[
        font=\small,
        box/.style={draw, rounded corners, align=center, minimum height=1cm,
                    text width=2.6cm, fill=black!3},
        arr/.style={-{Latex[length=2mm]}, thick}
    ]
        \node[box] (scan) {Scanner\\{\scriptsize tools $\cup$ assets}};
        \node[box, right=1.1cm of scan] (static) {Static scoring\\{\scriptsize 6 stages}};
        \node[box, right=1.1cm of static] (param) {Parameter\\rubric};
        \node[box, right=1.1cm of param] (call) {Call scoring\\{\scriptsize resolve + escalate}};
        \node[box, below=1.0cm of static, text width=6cm] (review)
            {Review: verify $\cdot$ judge $\cdot$ advise $\cdot$ critics};

        \draw[arr] (scan) -- (static);
        \draw[arr] (static) -- (param);
        \draw[arr] (param) -- (call);
        \draw[arr] (static) -- (review);
        \draw[arr] (call.south) |- (review.east);
    \end{tikzpicture}
    \caption{The McpRisk pipeline. The scanner and static-scoring stages run at
    design time; the parameter and call-scoring stages run at request time. The
    review subsystem audits outputs independently of the scorer.}
    \label{fig:pipeline}
\end{figure}

The static-scoring pipeline itself proceeds through six stages, each of which
prefers a local LLM working against an explicit rubric and degrades to a
deterministic fallback when the model is unavailable (except in \emph{strict}
mode, where an unreachable model raises an error rather than fabricating a
value): (0)~\emph{domain inference} of the server kind and profile;
(1)~\emph{tool impact} on a $1$--$3$ scale; (2)~\emph{asset sensitivity} on a
$1$--$5$ scale; (3)~\emph{blast radius} per (tool, asset) pair;
(4)~\emph{baselines} of expected behavior and \emph{risk-band} assignment; and
(5)~an independent \emph{judge} that re-derives every primitive from the same
domain profile and whose value overrides the proposer's on disagreement. Over
the seven evaluation servers this judge recorded $52$ overrides, and the
scan records a cross-check summary and provenance flag
(\texttt{llm-scan} vs.\ \texttt{offline-baseline}) for every cell.

\section{The Scoring Model}
\label{sec:framework:model}
\addcontentsline{tocheb}{section}{\protect\numberline{\secnumforhebrewtoc}{מודל הניקוד}}

At the conceptual level, McpRisk views the risk of an interaction as the
product of how much damage it can do and how likely that damage is:
\begin{equation}
\text{Risk} \;=\; \text{Impact} \times \text{Likelihood} \times
\text{Irreversibility}.
\label{eq:conceptual}
\end{equation}
In the implemented \emph{static} pipeline, likelihood is pinned to its upper
bound ($\text{Likelihood}=1.0$) because a design-time score must assume the
worst case, and irreversibility is folded into the tool-impact term. This
yields the concrete static score for a (tool, asset) cell:
\begin{equation}
\text{score} \;=\; \underbrace{s}_{\text{sensitivity}} \;\times\;
\underbrace{b}_{\text{blast radius}} \;\times\;
\underbrace{i}_{\text{tool impact}},
\qquad s \in [1,5],\; b \in [0,4],\; i \in [1,3],
\label{eq:static}
\end{equation}
so the score ranges over $[0, 60]$. The three primitives are defined as
follows.

\textbf{Sensitivity $s$ ($1$--$5$)} measures the confidentiality and value of
the asset being touched, from public data to regulated PII
(\Cref{tab:sensitivity}). \textbf{Blast radius $b$} measures how far a single
action propagates, from a single record (local) to the entire system
(systemic). \textbf{Tool impact $i$ ($1$--$3$)} measures reversibility: $1$ for
read-only tools that cannot change state, $2$ for tools that make a recoverable
change (write, edit, insert, move, post), and $3$ for irreversible or
access-changing tools (delete, drop, overwrite, execute, revoke, grant). This
last term is where the framework encodes the intuition that
\texttt{file\_write} is inherently riskier than \texttt{file\_read} and that
\texttt{delete\_file} is riskier still.

\begin{table}[H]
    \centering
    \caption{The asset-sensitivity scale ($s$).}
    \label{tab:sensitivity}
    \begin{tabular}{c l l}
        \hline
        \textbf{$s$} & \textbf{Asset class} & \textbf{Examples} \\
        \hline
        1 & Public / non-sensitive        & Public config, non-sensitive logs \\
        2 & Low-sensitivity infrastructure & Generic file store, API endpoints \\
        3 & Internal operational data     & Internal DB, system config, QA plans \\
        4 & Restricted business data      & Audit logs, source code, email \\
        5 & Regulated / PII               & Financial records, PII, private keys \\
        \hline
    \end{tabular}
\end{table}

The numeric score is discretized into four \textbf{risk bands} ---
\textit{low}, \textit{medium}, \textit{high}, \textit{critical} --- using the
thresholds and calibration rules of \Cref{tab:bands}. The calibration is
deliberately shaped like a \emph{risk pyramid}: only a cell that is irreversible
($i=3$), aimed at a crown-jewel asset ($s=5$), with at least departmental reach
($b \ge 3$) is labeled \textit{critical}, so that in practice only $1$--$2\%$ of
cells reach the top band. Two special rules encode confidentiality intuitions
that the raw product would miss: a broad read of highly sensitive data
($i=1, s\ge4, b\ge3$) is treated as mass exfiltration and raised to
\textit{high}, and any read of a crown-jewel asset ($i=1, s=5$) is never allowed
to fall below \textit{medium}. Read-only operations otherwise remain
\textit{low}, so that ordinary work is not gated.

\begin{table}[H]
    \centering
    \caption{Risk-band calibration. Thresholds on the raw score are
    $\text{medium}\ge 8$ and $\text{high}\ge 24$; the rules below refine them.}
    \label{tab:bands}
    \begin{tabular}{l l}
        \hline
        \textbf{Band} & \textbf{Condition} \\
        \hline
        critical & $i=3 \wedge s=5 \wedge b\ge 3$ \\
        high     & $(i=3 \wedge s\ge 4)\ \vee\ \text{score}\ge 24\ \vee\ (i=1 \wedge s\ge 4 \wedge b\ge 3)$ \\
        medium   & $\text{score}\ge 8\ \vee\ (i=1 \wedge s=5)$ \\
        low      & otherwise \\
        \hline
    \end{tabular}
\end{table}

\section{Atomic-Operation Taxonomy}
\label{sec:framework:atomic}
\addcontentsline{tocheb}{section}{\protect\numberline{\secnumforhebrewtoc}{טקסונומיית הפעולות האטומיות}}

To assign the tool-impact term consistently across heterogeneous servers, we
map every MCP tool to one or more of thirteen \emph{atomic operations}, each
carrying a fixed severity from $1$ (lightweight discovery) to $5$ (total
compromise), shown in \Cref{tab:atomic}. A tool's effective impact derives from
the highest-severity operation it can perform: a SQLite \texttt{write\_query}
tool that accepts free-form SQL, for example, is tagged with the full
worst-case set (\textsc{execute}, \textsc{delete}, \textsc{overwrite},
\textsc{schema\_modify}, \textsc{write}) because the argument alone determines
which it performs.

The mapping is produced by a \emph{rules-based} classifier that runs two
independent paths --- one over the server's README prose and one over the
structural signals of its tool-list schema (name prefixes, parameter shapes) ---
and surfaces any disagreement between them for review. Keeping this stage
rule-based rather than model-based makes the tool-to-severity mapping auditable
and fully reproducible.

\begin{table}[H]
    \centering
    \caption{The 13 atomic operations and their fixed severities.}
    \label{tab:atomic}
    \begin{tabular}{l c l}
        \hline
        \textbf{Operation} & \textbf{Severity} & \textbf{Rationale} \\
        \hline
        \textsc{execute}       & 5 & Arbitrary code/command execution; total compromise \\
        \textsc{delete}        & 5 & Irreversible data destruction \\
        \textsc{overwrite}     & 4 & Corrupts existing data in full \\
        \textsc{schema\_modify} & 4 & Structural damage; cascading failures \\
        \textsc{broadcast}     & 4 & Sends data externally; exfiltration + impersonation \\
        \textsc{write}         & 3 & Injects new data; malicious persistence \\
        \textsc{modify}        & 3 & Tampers with existing data; integrity violation \\
        \textsc{move}          & 3 & Relocates resources; can hide or expose data \\
        \textsc{create}        & 3 & New persistent resource; staging for writes \\
        \textsc{read}          & 2 & Accesses content; confidentiality breach \\
        \textsc{search}        & 2 & Targeted discovery of sensitive data \\
        \textsc{metadata}      & 1 & Attributes without content; structural recon \\
        \textsc{list}          & 1 & Enumerates resources; existence discovery \\
        \hline
    \end{tabular}
\end{table}

\section{Parameter (Request-Time) Scoring}
\label{sec:framework:params}
\addcontentsline{tocheb}{section}{\protect\numberline{\secnumforhebrewtoc}{ניקוד ברמת הפרמטר (בזמן ריצה)}}

The static cell score is an upper bound over all possible invocations of a tool.
To turn it into a request-time score, McpRisk adds a \emph{parameter} layer that
reads the actual arguments of a call. For each tool, the LLM derives a rubric
that identifies the magnitude-bearing parameters, and for each such parameter a
\texttt{base\_rank} (its intrinsic sensitivity), an extraction rule (whether the
magnitude is a number, a list length, a parsed \texttt{LIMIT}, or a boolean),
and a set of ascending numeric \texttt{cutoffs}. A concrete value is banded by
the highest cutoff it meets. A crucial rule is that the \emph{absence} of a
\texttt{LIMIT} clause is treated as unbounded ($\infty$) and therefore
worst-case, so that a \texttt{SELECT} or \texttt{DELETE} touching every row is
scored as maximally broad.

Two bands are then combined arithmetically, where
$\text{low}=1,\ \text{medium}=2,\ \text{high}=3,\ \text{critical}=4$. The
\emph{parameter risk} is the average of the parameter's base rank and its value
band, rounded half-up:
\begin{equation}
\text{param\_risk} \;=\; \left\lfloor \tfrac{\text{base\_rank} +
\text{value\_band}}{2} + 0.5 \right\rfloor .
\label{eq:paramavg}
\end{equation}
The final band of the call escalates the inherent (tool, asset) band by the
parameter risk, taking the more severe of the two so that a risky argument can
only \emph{raise} the band, never lower it:
\begin{equation}
\text{final\_band} \;=\; \max\!\big(\text{cell\_band},\ \text{param\_risk}\big).
\label{eq:escalate}
\end{equation}
For example, a \texttt{create\_event} call whose \texttt{attendees} parameter has
base rank \textit{medium} and $50$ recipients (value band \textit{critical})
yields $\text{param\_risk}=\lfloor(2+4)/2+0.5\rfloor=3$; combined with a
\textit{high} inherent cell this escalates the call to \textit{critical}. When
several parameters apply, the most severe governs.

\section{Operationalizing Likelihood: Misuse Scoring of Call Logs}
\label{sec:framework:misuse}
\addcontentsline{tocheb}{section}{\protect\numberline{\secnumforhebrewtoc}{ניקוד שימוש לרעה של יומני קריאות}}

The parameter layer captures the \emph{magnitude} of a request, but the
conceptual model of \Cref{eq:conceptual} also has a \emph{likelihood} term. To
score a captured call log end to end, McpRisk instantiates that term with a
concrete formula, calibrated against the canonical lookup tables used to build
the project's hand-labeled heatmaps:
\begin{align}
\text{AssetImpact} &= \min(\text{Sensitivity}\times\text{Blast},\ 25),
\label{eq:assetimpact}\\
\text{Likelihood}  &= \text{BaseDev}\times\text{Anomaly}\times\text{Exposure},
\label{eq:likelihood}\\
\text{Misuse}      &= \max\!\big(\text{AssetImpact}\times\text{Likelihood}
                       \times\text{ToolImpact},\ 1.0\big),
\label{eq:misuse}
\end{align}
where $\text{ToolImpact}\in\{1,2,3\}$ is the reversibility multiplier of
\Cref{sec:framework:model}. The three likelihood sub-factors follow the
definitions of the conceptual model: \emph{baseline deviation} (how unusual the
call is relative to intended use), \emph{access anomaly} (time/volume
outlierness), and \emph{exposure} (how reachable the destination makes the
asset). Crucially, the first two require a history of prior calls that a server
gating a request in isolation does not have; consistent with the framework's
open question about how much runtime log context is available, they are set to
conservative constants ($\text{BaseDev}=0.9$, a high prior that any logged call
may deviate; $\text{Anomaly}=0.3$, a low value reflecting the absence of
observable access patterns), while \emph{exposure} is read from a lookup table
keyed on the destination directory (from $0.3$ for public areas to $0.9$ for
security-evidence stores). The floor of $1.0$ ensures every logged call registers
at least baseline activity. Because the context multiplier of the original
design was removed, the band thresholds are rescaled accordingly to
$\textit{critical}\ge 27$, $\textit{high}\in[13,26]$, $\textit{medium}\in[7,12]$,
and $\textit{low}<7$.

Five judgment rules close gaps the lookup tables leave open, and encode the
threat intuitions a naive product would miss. (i)~For cross-path tools
(move/copy/rename), sensitivity is taken from the \emph{source} and exposure from
the \emph{destination}; if a sensitive source ($\text{Sensitivity}\ge 4$) is
moved to a low-exposure public area, exposure is forced to $1.0$ --- the asset is
now reachable to anyone, and renaming its extension does not lower its
sensitivity. (ii)~A path-traversal attempt (e.g.\ \texttt{../../etc/passwd}) is
scored as a system asset ($\text{Sensitivity}=5$, $\text{Exposure}=0.9$) even if
it was blocked, because the probe itself is the signal. (iii)~For multi-path
calls, the worst case across paths governs. (iv)~Errors and denied calls still
score, since the attempt is a signal. (v)~Asset-name overrides raise sensitivity
where the extension hides it --- an \texttt{audit\_log.txt} or a file matching
\texttt{*credential*} is scored at sensitivity $5$ regardless of its
\texttt{.txt} suffix. These rules are exactly the cases where the static cell
score under-weights a real threat, and folding them into the request-time score
is how the dynamic mode adds value over the static upper bound. A richer
contextual scorer that learns baseline deviation and access anomaly from real
runtime logs remains future work, as we discuss in
\Cref{chapter:discussion_and_conclusions}.

\section{Implementation and Reproducibility}
\label{sec:framework:implementation}
\addcontentsline{tocheb}{section}{\protect\numberline{\secnumforhebrewtoc}{מימוש ושחזוריות}}

McpRisk is implemented in Python~$3.12$ as the package
\texttt{mcp\_security}, with the stages of \Cref{fig:pipeline} realized as the
sub-packages \texttt{scanner}, \texttt{static\_scoring},
\texttt{param\_scoring}, \texttt{call\_scoring}, \texttt{atomic\_ops}, and
\texttt{review}. All model-backed stages call a single local model ---
Qwen2.5:32B served through Ollama --- with no cloud fallback by design, so that
no server metadata leaves the machine. Decoding is made deterministic and
reproducible by fixing greedy decoding parameters (temperature $0$, seed $0$,
top-$p$ $1.0$). Because the same model both proposes and (through the judge
stage) cross-checks scores, and because it also underlies part of the
evaluation, we treat model-induced circularity as a first-class limitation and
address it in \Cref{chapter:experiments} by validating against \emph{external}
human and framework oracles.
